\section{Voltage inequality as a quadratic form}

Substituting \cref{eq:affine-map} into \cref{eq:voltage-ball} gives
\begin{align}
0&\geq \norm{\Apsm\current+\vect r}^2-\Umax^2\\
 &=\current\trans\Apsm\trans\Apsm\current
   +2\vect r\trans\Apsm\current+\vect r\trans\vect r-\Umax^2.
\end{align}
Thus
\begin{equation}
 \boxed{\current\trans\Hpsm\current+2\vect b\trans\current+c\leq0}
 \label{eq:psm-general-quadratic}
\end{equation}
with
\begin{align}
\Hpsm&=\Apsm\trans\Apsm
=\begin{bmatrix}
\Rs^2+\omegae^2\Ld^2&\Rs\omegae(\Ld-\Lq)\\
\Rs\omegae(\Ld-\Lq)&\Rs^2+\omegae^2\Lq^2
\end{bmatrix},\label{eq:Hpsm}\\
\vect b&=\Apsm\trans\vect r
=\begin{bmatrix}\omegae^2\Ld\psif\\\Rs\omegae\psif\end{bmatrix},
\label{eq:bpsm}\\
c&=\omegae^2\psif^2-\Umax^2.\label{eq:cpsm}
\end{align}
Every term now has a traceable origin.  The diagonal entries of $\Hpsm$ are
squared voltage sensitivities to $i_d$ and $i_q$.  The off-diagonal term comes
from resistance acting together with saliency.  The linear vector couples the
PM-generated voltage to current-generated voltage.  The constant compares
open-circuit back EMF with available voltage.

The matrix $\Hpsm$ is symmetric by construction.  Moreover,
\begin{equation}
 \vect y\trans\Hpsm\vect y=\norm{\Apsm\vect y}^2\geq0,
\end{equation}
so it is positive semidefinite without any appeal to an expected ellipse.  Its
determinant and trace are
\begin{align}
\det\Hpsm&=(\det\Apsm)^2
=(\Rs^2+\omegae^2\Ld\Lq)^2,\label{eq:detH}\\
\tr\Hpsm&=2\Rs^2+\omegae^2(\Ld^2+\Lq^2)>0
\end{align}
except at the simultaneous ideal limit $R_s=\omega_e=0$.  Therefore, whenever
$\Rs^2+\omegae^2\Ld\Lq>0$, both eigenvalues are positive and $\Hpsm$ is
positive definite.

Why does the determinant enter here?  For a linear map in the plane,
$\abs{\det\Apsm}$ is the factor by which oriented area is scaled.  A nonzero
determinant says that no area has collapsed to a line and that every voltage
increment has one current increment as its inverse.  If the determinant is
zero, at least one current direction is collapsed: that direction lies in the
null space and produces no voltage increment.  The identity
$\det(\Apsm\trans\Apsm)=(\det\Apsm)^2$ therefore translates the geometric
area test into a positivity test for the quadratic form.  It also explains
why rank loss, a zero eigenvalue, a null direction, and a degenerate or open
level set are different descriptions of the same collapse.

\physinterpretation No nonzero current displacement is invisible to both
voltage components when $\Apsm$ is invertible.

\engineeringconsequence The closedness of the eventual current-space boundary
comes from full column rank of the current-to-voltage map, not from the label
``synchronous machine.''

\paragraph{Synthesis.}
Algebra gives a Gram matrix $\Hpsm=\Apsm\trans\Apsm$; geometry says every
nonzero direction has positive quadratic growth; physics says every current
direction changes at least one voltage component; engineering says a finite
voltage disk pulls back to a bounded current set.  This chain fails precisely
when the map loses rank.
